\begin{lemma}\label{lem:balls}
Let $(\Sigma,d)$ be $\mathbb R^3$, hyperbolic space $H^3$, or the unit sphere
$S^3$ with $r_2\le\pi$. Then the closed ball $\bar B(X,r_1)$ lies in the open
ball $B(Y,r_2)$ if and only if $d(X,Y)+r_1<r_2$.
\end{lemma}

\begin{proof}
($\Leftarrow$) is the triangle inequality. ($\Rightarrow$) Let $r_1\ge0$ and suppose
$d(X,Y)+r_1\ge r_2$. On $\mathbb R^3$ and $H^3$, where geodesics extend
indefinitely and minimize, or on $S^3$ when $d(X,Y)+r_1\le\pi$, follow a minimizing geodesic from $Y$ through $X$ (any geodesic from
$X$ if $X=Y$) a further distance $r_1$; the endpoint lies in $\bar B(X,r_1)$ at
distance $d(X,Y)+r_1\ge r_2$ from $Y$. If $d(X,Y)+r_1>\pi$ on $S^3$, the
antipode of $Y$ is at distance $\pi-d(X,Y)<r_1$ from $X$ and $\pi\ge r_2$ from
$Y$.
\end{proof}

\begin{proposition}\label{prop:conformal}
Let $M$ be conformal to $(\eta_-,\eta_+)\times\Sigma$ with metric
$-d\eta^2+d\Sigma^2$, where $\Sigma$ is $\mathbb R^3$, $H^3$, or the unit $S^3$
with $\eta_+-\eta_-\le\pi$. Then $\Det(b)=\Is(b)$ for every $b$.
\end{proposition}

\begin{proof}
Use $t=\eta$. For $x=(\eta_x,X)$, $b=(\eta_b,Y)$: $x\ll b$ iff
$d(X,Y)<\eta_b-\eta_x$ (on $S^3$, $\eta_b-\eta_x<\pi$). For $\eta_-<c<\eta_x$,
$J^-(x)\cap S_c=\bar B(X,\eta_x-c)$ and $\Is(b)\cap S_c=B(Y,\eta_b-c)$, with
$\eta_b-c<\pi$ on $S^3$. By Lemma~\ref{lem:balls} the inclusion holds iff
$d(X,Y)<\eta_b-\eta_x$, i.e.\ iff $x\ll b$. Apply Theorem~\ref{thm:char}.
\end{proof}

This covers Minkowski spacetime, spatially flat and open FLRW universes with
a big bang (conformal to $\eta>0$), the closed radiation universe
($a\propto\sin\eta$, $\eta\in(0,\pi)$) and global de Sitter spacetime ($\Sigma=S^3$,
$\eta\in(-\pi/2,\pi/2)$) \citep{HE1973}. In none of them is anything outside an
event's chronological past guaranteed from it. In Minkowski spacetime and the
spatially flat universes this bound is attained by a natural field.

\begin{proposition}\label{prop:sharp}
Let $M=(\eta_-,\eta_+)\times\mathbb R^3$ with $g=a(\eta)^2(-d\eta^2+d\vec X^2)$,
$a>0$ smooth, and let $\mathcal L$ be the smooth solutions of the conformally
coupled scalar equation $\Box_g\phi-\tfrac16R\phi=0$. Then
$F_{\mathcal L}(b)=J^-(b)$ for every $b$.
\end{proposition}

\begin{proof}
$J^-(b)\subseteq F_{\mathcal L}(b)$ by continuity, as in
Remark~\ref{rem:laws}. Conversely, $\phi$ is a solution iff $\psi=a\phi$
solves the flat wave equation $\partial_\eta^2\psi=\Delta\psi$ on $M$
\citep[app.~D]{Wald1984}. Let $b=(\eta_b,\vec Y)$ and $x=(\eta_x,\vec X)\notin
J^-(b)$, so $\eta_x-\eta_b+|\vec X-\vec Y|>0$, since the causal curves of $M$
are those of the conformal Minkowski metric and the slab is causally convex in
Minkowski spacetime. Choose a unit vector $\hat k$
with $\hat k\cdot(\vec X-\vec Y)=-|\vec X-\vec Y|$, a smooth $f$ vanishing on
$(-\infty,0]$ and positive on $(0,\infty)$, and put $\psi=f(u)$ with
$u=\eta-\eta_b-\hat k\cdot(\vec Z-\vec Y)$ at $(\eta,\vec Z)$. This solves the
wave equation. On $\Is(b)$, $\eta-\eta_b<-|\vec Z-\vec Y|$, so $u<0$ and
$\psi=0$; at $x$, $u>0$ and $\psi\neq0$. So $\psi/a$ and $0$ are solutions that
agree on $\Is(b)$ but not at $x$.
\end{proof}


\begin{proposition}\label{prop:dust}
Let $M$ be conformal to $(\eta_-,\eta_+)\times S^3$ with $\eta_+-\eta_->\pi$.
Then $\Det(b)=\Is(b)$ if $\eta_b\le\eta_-+\pi$, and $\Det(b)=M$ if
$\eta_b>\eta_-+\pi$. In the Einstein static universe ($\eta_\pm=\pm\infty$),
$\Det(b)=M$ for every $b$.
\end{proposition}

\begin{proof}
If $\eta_b>\eta_-+\pi$, take $c\in(\eta_-,\eta_b-\pi)$: $\Is(b)\cap S_c$ is an open
ball of radius greater than $\pi$, hence all of $S_c$, and
Corollary~\ref{cor:cauchy} applies. If $\eta_b\le\eta_-+\pi$, every $c>\eta_-$
has $\eta_b-c<\pi$, and the proof of Proposition~\ref{prop:conformal} applies.
\end{proof}

For the closed dust FLRW model, $a(\eta)=\tfrac12a_m(1-\cos\eta)$ with
$\eta\in(0,2\pi)$, the threshold is $\eta=\pi$, the moment of maximum
expansion. That the particle horizon of the closed dust model disappears at
maximum expansion, so that an observer can then in principle see the whole
universe, is known \citep{GronJohannesen2011,BolotinTanatarov2013}; the
Einstein static case is Geroch's \citep[Fig.~5]{Geroch1977}. What
Proposition~\ref{prop:dust} adds is the consequence for determination: with
Corollary~\ref{cor:cauchy}, once the particle horizon has vanished
($\eta_b>\pi$), everything is guaranteed from $b$, so that every law with the
domain-of-dependence property fixes everything from what is real from $b$.
This fits Manchak's theorem that if some event of $b$'s timelike future lies
in Geroch's domain of prediction of $b$, the spacetime admits a compact
spacelike slice \citep[p.~319]{Manchak2008}. That theorem concerns the
geometric notion. The flow law of Remark~\ref{rem:laws} shows that it has no
analogue for what a particular law fixes: in Minkowski spacetime, which has no
compact slice, that law fixes all of $b$'s timelike future. Manchak also
records the view that ``Maxwell's equations constrain some spacetimes in such
a way as to allow for prediction'' \citep[p.~321 and n.~15]{Manchak2008}, so
the law-relative notion has antecedents. Earman's discussion explicitly
includes the Cauchy-surface exception and an electromagnetic-constraint
example citing Ellis and Sciama \citep[pp.~125--130]{Earman1995}. We have checked
these selected original Earman pages; Ellis and Sciama (1972, pp.~41--42)
remain unread here and enter this comparison through Earman's discussion.
The model-specific characterization below may overlap with those sources;
we make no priority claim for it.

Before maximum expansion a particular field can fix more than is guaranteed.
Light leaving a point of the closed model refocuses at its antipode after
conformal time $\pi$, and the rescaled conformal scalar $\psi=a\phi$ returns
there with its sign reversed.

\begin{proposition}\label{prop:echo}
Let $M$ be the closed dust model, $g=a(\eta)^2(-d\eta^2+d\Omega_3^2)$ with
$a=\tfrac12a_m(1-\cos\eta)$, $\eta\in(0,2\pi)$, and $\mathcal L$ as in
Proposition~\ref{prop:sharp}. Let $A(\eta,X)=(\eta+\pi,-X)$, where $-X$ is the
antipode of $X$. Then $J^-(b)\cup\bigl(A(J^-(b))\cap M\bigr)\subseteq
F_{\mathcal L}(b)$. The set $A(J^-(b))\cap M$ is disjoint from $\Isp(b)$, and if
$\eta_b\le\pi$ it is nonempty and disjoint from $\Det(b)$.
\end{proposition}

\begin{proof}
Here $\psi=a\phi$ solves the conformally invariant equation of the Einstein
static universe, $\partial_\eta^2\psi=\Delta\psi-\psi$ on the unit $S^3$
\citep[app.~D]{Wald1984}, and by global well-posedness extends uniquely to
$\mathbb R\times S^3$.
Expanding in spherical harmonics of degree $n$, which have frequency $n+1$
and parity $(-1)^n$ under the antipodal map, each mode, and so every smooth
solution, satisfies $\psi(\eta+\pi,-X)=-\psi(\eta,X)$. Two solutions that agree on $\Is(b)$ agree, by continuity, on $J^-(b)$, and
hence on $A(J^-(b))\cap M$. If $(\eta,X)\in J^-(b)$ and $\eta+\pi<2\pi$, then
$d(X,Y)\le\eta_b-\eta$, so $d(-X,Y)=\pi-d(X,Y)\ge(\eta+\pi)-\eta_b$;
thus $A(\eta,X)\notin\Isp(b)$. If $\eta_b\le\pi$, then
$\Det(b)=\Is(b)$ by Proposition~\ref{prop:dust}, which lies in
$\{\eta<\eta_b\}$, while $A(J^-(b))$ lies in $\{\eta>\pi\}$ and meets $M$, for
instance at $A(\eta_b/2,Y)$ when $b=(\eta_b,Y)$.
\end{proof}

For $\eta_b\le\pi$, the echo events are unreal from $b$ and fixed from $b$,
but they lie beside $b$ or on its future light cone, not in its timelike
future. The field fixes
more than the echo, including part of $b$'s timelike future. Let
$\Diamond_b=\{(\eta,X)\in\mathbb R\times S^3:d(X,Y)+|\eta|<\eta_b\}$, the double
cone in the Einstein static universe between $b$ and its mirror image
$(-\eta_b,Y)$, and $U_b=\Diamond_b\cup A\Diamond_b\cup A^2\Diamond_b$.

\begin{theorem}\label{thm:diamonds}
In the setting of Proposition~\ref{prop:echo}, let $0<\eta_b<\pi$. Then
$F_{\mathcal L}(b)=M\cap\overline{U_b}$. In
particular $F_{\mathcal L}(b)$ contains the nonempty set
$\{(\eta,X)\in M:d(X,Y)\le\eta+\eta_b-2\pi\}$, which lies in $\Isp(b)$.
\end{theorem}

\begin{proof}
For data $(f,g)$ at $\eta=0$, the Einstein static solution is
$\psi(\eta,X)=\partial_\eta[\sin\eta\,M_f(X,\eta)]+\sin\eta\,M_g(X,\eta)$, where
$M_h(X,r)$ is the mean of $h$ over the geodesic sphere of radius $r$ about
$X$; one checks this on each mode, using that the zonal function of degree $n$
on $S^3$ is $\sin((n+1)r)/((n+1)\sin r)$. So $\psi(x)$ depends on the data only
near the sphere $\Sigma_x=S(X,\eta_x)$, which is the point $-X$ when $\eta_x=\pi$.

\emph{Sufficiency.} Solutions agreeing on $\Is(b)$ agree on the ball
$B(Y,\eta_b-\epsilon)$ in the slice $\eta=\epsilon$, hence on its domain of
dependence in the Einstein static universe; the union over $\epsilon>0$ is
$\Diamond_b$. The identities $\psi\circ A=-\psi$ and $\psi\circ A^2=\psi$ carry
the agreement to $A\Diamond_b$ and $A^2\Diamond_b$.
Moreover, $F_{\mathcal L}(b)$ is relatively closed in $M$: it is the
intersection of the zero sets of all smooth difference solutions vanishing
on $\Is(b)$. Since $M$ is open in the Einstein cylinder, taking the relative
closure gives $M\cap\overline{U_b}\subseteq F_{\mathcal L}(b)$.

\emph{Necessity.} Let $x\in M\setminus\overline{U_b}$. Replacing $x$ by
$A^{-1}x$ if $\eta_x>\pi$, assume $\eta_x\le\pi$. Let $C=\bar B(Y,\eta_b)$.
Data supported off $C$ give a solution vanishing on $\Is(b)$: if $z\in\Is(b)$
lay in its support, then $w\le z\ll b$ for some $w$ in the support of the data,
so $d(w,Y)<\eta_b$. If $\eta_x=\pi$, then $d(-X,Y)>\eta_b$, and data
$(f,0)$ with $f$ a bump at $-X$ off $C$ give $\psi(x)=-f(-X)\neq0$. If
$\eta_x<\pi$, then $\Sigma_x\not\subseteq C$: otherwise the connected set
$S^3\setminus C$ lies in one of the two open balls bounded by $\Sigma_x$, so
the other closed ball lies in $C$, which by the argument of
Lemma~\ref{lem:balls} means
$d(X,Y)+\eta_x\le\eta_b$ or $d(-X,Y)+\pi-\eta_x\le\eta_b$, that is,
$x\in\overline{\Diamond_b}\cup\overline{A\Diamond_b}$. Take $p\in\Sigma_x\setminus C$
and data $(0,g)$ with $g\ge0$ supported near $p$ off $C$ and $g(p)>0$; then
$\psi(x)=\sin\eta_x\,M_g(X,\eta_x)>0$.

Finally, if $d(X,Y)\le\eta+\eta_b-2\pi$, then $d(X,Y)<\eta-\eta_b$ because
$\eta_b<\pi$, so $(\eta,X)\in\Isp(b)$; the set contains $(\eta,Y)$ for
$\eta\ge2\pi-\eta_b$.
\end{proof}

On $b$'s comoving geodesic $X=Y$, the field is fixed from $b$ exactly for
$\eta\le\eta_b$ and for $\eta\ge2\pi-\eta_b$, and open in between: the end is
fixed before the middle. The reason is that the past light cone of an event
just before the crunch wraps around the universe and refocuses near $Y$ at the
big bang, inside what is real from $b$. The same sphere argument shows that,
up to boundary, $M\cap U_b$ is the set
of events all of whose past-inextendible null geodesics meet $\Is(b)$, so for
this field, which propagates on light rays alone, the null geodesics decide
what is fixed. The same argument applies to global de Sitter spacetime,
conformal to the slab $(-\pi/2,\pi/2)$, where the field is fixed from $b$ on the
antipodal comoving geodesic for all $\eta>-\eta_b$, although no signal from that
geodesic ever reaches the comoving geodesic of $b$. So the geometric threshold
at maximum expansion marks when everything is guaranteed, not when this field
first fixes anything beyond the past.

For a law whose fixed set itself changes at the threshold, take the comoving
flow $u=a^{-1}\partial_\eta$ of the closed dust model with the equations of
Remark~\ref{rem:laws}; here $u$ is the dust's own four-velocity, so the law
adds no structure the model lacks. The comoving line through $X$ meets $\Is(b)$ iff
$d(X,Y)<\eta_b$, so, by continuity in $X$ and by choosing data supported off
the closed ball, $F_{\mathcal L}(b)=(0,2\pi)\times\bar B(Y,\eta_b)$. This is a
proper subset of $M$, containing events in $b$'s timelike future, if
$\eta_b<\pi$, and all of $M$ if $\eta_b\ge\pi$. Under this law the answer to
``is the future fixed from here?'' differs between events before and after
maximum expansion.

Theorem~\ref{thm:massive} shows that on closed dust a local law can fix exactly
$\Det(b)\cup J^-(b)$ for every $\tau\neq\pi$. Whether some local law does so at
$\tau=\pi$, or in other closed models, we leave open.
